\chapter{Lecture 10: Copy Constructor, Destructor and Copy Assignment}
\begin{center}
\textit{February 6, 2018}
\end{center}

\section{Recap: Member Initialization List}
\begin{itemize}
  \item must be used to initialize \code{const} and reference fields;
  \item can be used for all fields;
  \item initialization occurs in declaration order;
  \item can be more efficient than initializing in the ctor body.
\end{itemize}
\begin{lstlisting}[style=cpp]
Student::Student(int id, int assns, int mt, int final)
    : id{id}, assns{assns}, mt{mt}, final{final} {}
\end{lstlisting}
The name outside the \code{\{\}} is the field; the expression inside the
\code{\{\}} is the value.

\section{The Big Five}
\begin{lstlisting}[style=cpp]
Student billy{60, 70, 80};
Student bobby{billy};      // or: Student bobby = billy;
\end{lstlisting}
Objects are copied by calling the \textbf{copy constructor}, and you get one
for free. Every class comes with:
\begin{enumerate}
  \item a default constructor (only if you write no ctor),
  \item a copy constructor,
  \item a destructor,
  \item a copy assignment operator,
  \item a move constructor (C++11),
  \item a move assignment operator (C++11).
\end{enumerate}
Items 2--6 are called the \textbf{Big Five}.

\begin{quote}
\textit{``There will be a question about them, and I have already made that
question an hour ago.''} --- N.~Naeem
\end{quote}

\section{The Copy Constructor}
The free copy ctor does a field-for-field copy. For \code{Student} it is
equivalent to:
\begin{lstlisting}[style=cpp]
Student::Student(const Student &other)
    : assns{other.assns}, mt{other.mt}, final{other.final} {}
\end{lstlisting}
\begin{quote}
\textit{``Beggars can't be choosers.''} --- N.~Naeem
\end{quote}
Sometimes the free copy ctor does not do the right thing:
\begin{lstlisting}[style=cpp]
// node.h (include guard assumed)
struct Node {
    int data;
    Node *next;
    Node(int data, Node *next);
    Node(const Node &other);
};

// node.cc
Node::Node(int data, Node *next) : data{data}, next{next} {}
Node::Node(const Node &other) : data{other.data}, next{other.next} {}  // = free one

// main.cc
Node *np = new Node{1, new Node{2, new Node{3, nullptr}}};
Node m{*np};               // copy ctor; m is on the stack
Node *m2 = new Node{*np};  // copy ctor; *m2 is on the heap
\end{lstlisting}

\begin{center}
\begin{tikzpicture}[box/.style={draw, minimum width=13mm, minimum height=7mm},
  >={Stealth}]
\node[box] (np) at (0,0) {\code{np}};
\node[box] (a) at (2.5,0) {1 \,$\bullet$};
\node[box] (b) at (5,0) {2 \,$\bullet$};
\node[box] (c) at (7.5,0) {3 \,/};
\node[box] (m) at (2.5,-1.5) {1 \,$\bullet$};
\node[left=2mm of m] {\code{m}};
\draw[->] (np) -- (a);
\draw[->] (a) -- (b);
\draw[->] (b) -- (c);
\draw[->] (m.east) -- (b.south west);
\end{tikzpicture}
\end{center}

This is not really a copy of the list: \code{m} shares nodes 2 and 3 with the
original. This is called a \textbf{shallow copy}. For a pointer field we usually
want to copy the object the pointer points to: a \textbf{deep copy}.

\begin{lstlisting}[style=cpp]
Node::Node(const Node &other)
    : data{other.data},
      next{other.next ? new Node{*other.next} : nullptr} {}
\end{lstlisting}
\code{new Node\{*other.next\}} is a recursive call to the copy ctor, which
copies the rest of the list. The check is needed because \code{other.next} could
be \code{nullptr}, and dereferencing it would crash.

\subsection{When Is the Copy Constructor Called?}
\begin{enumerate}
  \item when an object is initialized from another object of the same type;
  \item when an object is passed by value;
  \item when an object is returned by value.
\end{enumerate}
(2 and 3 may be optimized away or replaced by a move; see Lecture~11--12.)

\begin{important}[Why the parameter is a reference]
Because of (2), the copy ctor's parameter must be passed by reference: if it were
\code{Node(Node other)}, passing the argument would need a copy, which calls the
copy ctor, which needs a copy, \dots{} infinite recursion. The language forces
the reference.
\end{important}

\section{Single-Parameter Constructors}
\begin{lstlisting}[style=cpp]
Node::Node(int data) : data{data}, next{nullptr} {}

Node n{4};
Node n2 = 4;        // also compiles!
void foo(Node n);
foo(4);             // also compiles: 4 is converted to a Node
\end{lstlisting}
One-parameter ctors create \emph{implicit} (automatic) conversions. This is
exactly how \code{string s = "hello";} works: \code{std::string} has a
one-parameter ctor taking a \code{const char*}. But for \code{Node}, silently
turning an \code{int} into a Node is surprising and can hide bugs. To disable
the conversion, use \code{explicit}:
\begin{lstlisting}[style=cpp]
struct Node {
    explicit Node(int data);
    ...
};
Node n{4};    // OK
Node n2 = 4;  // error
foo(4);       // error
foo(Node{4}); // OK
\end{lstlisting}

\section{The Destructor}
When an object is destroyed, a method called the \textbf{destructor} (dtor)
runs. Objects are destroyed when
\begin{enumerate}
  \item a stack-allocated object goes out of scope, or
  \item a heap object is explicitly \code{delete}d.
\end{enumerate}

\begin{important}[Steps of object destruction]
\begin{enumerate}
  \item The dtor body runs.
  \item Fields are destroyed: dtors are called for fields that are objects (in
        reverse declaration order).
  \item The space is reclaimed.
\end{enumerate}
\end{important}

A class has exactly one dtor, named \code{\textasciitilde ClassName()}, with no
parameters and no return type. The free dtor has an empty body (and still does
step 2). But that might not be what we want:
\begin{lstlisting}[style=cpp]
Node *np = new Node{1, new Node{2, new Node{3, nullptr}}};
delete np;    // calls the dtor for the node with value 1
\end{lstlisting}
\code{next} is a pointer, not an object, so step 2 does nothing for it: only the
first node is freed, and nodes 2 and 3 are leaked. Let's write our own dtor:
\begin{lstlisting}[style=cpp]
Node::~Node() { delete next; }
\end{lstlisting}
This recursively deletes the rest of the list; the recursion ends at the last
node because \code{delete nullptr} does nothing.

\section{The Copy Assignment Operator}
\begin{lstlisting}[style=cpp]
Student billy{60, 70, 80};
Student bobby{billy};   // copy ctor: bobby is being created
Student jane;
jane = billy;           // copy assignment: jane already exists
                        // same as jane.operator=(billy);
\end{lstlisting}
The free \code{operator=} does a field-for-field copy. Once a class manages heap
memory, we need to write our own. Let's implement \code{operator=} for
\code{Node}.

First, what should it return? Assignment associates to the right:
\code{a = b = c = 0;} assigns \code{c = 0}, which produces \code{c}; then
\code{b = c} produces \code{b}; then \code{a = b}. For objects,
\begin{lstlisting}[style=cpp]
n1 = n2 = n3;   // n1.operator=(n2.operator=(n3));
\end{lstlisting}
So \code{operator=} must return the object (by reference, to avoid a copy) to
allow cascaded assignment.

\begin{lstlisting}[style=cpp]
Node &Node::operator=(const Node &other) {
    if (this == &other) return *this;   // self-assignment check: see below
    data = other.data;
    delete next;     // important: next may already own heap memory
    next = other.next ? new Node{*other.next} : nullptr;
    return *this;
}
\end{lstlisting}
The MIL cannot be used here: it is only available for constructors (the object
already exists).

\subsection{Why Check for Self-Assignment?}
\begin{lstlisting}[style=cpp]
Node n{1, new Node{2, nullptr}};
n = n;       // n.operator=(n): unlikely to be written directly,
             // but *p = *q where p == q happens easily
\end{lstlisting}
Without the check, \code{other.next} is the same pointer as \code{next}. We
\code{delete next}, which makes \code{other.next} a dangling pointer, and then
\code{new Node\{*other.next\}} copies freed memory. Solution: check for
self-assignment first.

\paragraph{Q.} What happens if \code{new} fails (no memory)?

Then \code{next} has already been deleted, but has not been given a new value:
the object is left holding a dangling pointer. Next lecture fixes this.

\section*{Summary}
\begin{itemize}
  \item The Big Five: copy ctor, copy assignment, dtor, move ctor, move
        assignment. The free versions copy/move field by field.
  \item With pointer fields, the free copy is shallow; write a deep copy.
  \item The copy ctor takes \code{const T\&} and is used for initialization,
        pass-by-value and return-by-value.
  \item Mark one-parameter ctors \code{explicit} unless you want implicit
        conversion.
  \item Dtor: body, then fields, then space. Write one to free owned heap memory.
  \item \code{operator=} returns \code{*this} by reference, must free the old
        resources, and must handle self-assignment.
\end{itemize}
